\documentclass[aspectratio=169]{beamer}
\mode<presentation>
\usetheme{default}
\usepackage{tikz}

% to keep this somewhat sane
\newcommand{\homebuttons}{
    % https://tex.stackexchange.com/questions/187993/how-to-position-beamer-button-exactly
    \begin{tikzpicture}[remember picture, overlay]
        \node[shift={(-0.5cm, -0.5cm)}]() at (current page.north east){
            \hyperlink{Usage}{\beamerbutton{Help}}
        };
        \node[shift={(-1.0cm, 0.5cm)}]() at (current page.south east){
            \hyperlink{AHome}{\beamerbutton{A}}
        };
        \node[shift={(-0.5cm, 0.5cm)}]() at (current page.south east){
            \hyperlink{BHome}{\beamerbutton{B}}
        };
    \end{tikzpicture}
}
\newcommand{\usebackground}[2][1.0]{
    % https://tex.stackexchange.com/questions/7916/how-to-insert-a-background-image-in-a-beamer-frame
    % https://tex.stackexchange.com/questions/475375/absolute-positioning-of-images-in-beamers-background-canvas?noredirect=1&lq=1
    \usebackgroundtemplate{
        \vbox to \paperheight{\vfil\hbox to \paperwidth{\hfil
                \tikz\node[opacity=#1, inner sep=0pt]{\includegraphics[width=\paperwidth, height=\paperheight, keepaspectratio]{#2}};
        \hfil}\vfil}
    }
}

\begin{document}
\beamertemplatenavigationsymbolsempty
\setbeamercolor{button}{bg=blue, fg=white}
\title[SCP]{Short Creative Project}
\subtitle{HISTART 392}
\author{Johnson He (hejohns)}
\date{2021-03-08}
\begin{frame}
    \titlepage
    \homebuttons
\end{frame}
\begin{frame}
    \frametitle{About}
    \begin{itemize}
        \item Made using beamer (and a lot of hyperlinks and a little TikZ)
        \item That said, this is NOT a slide show
        \item See Usage/Help
        \item Sources at the very end
    \end{itemize}
    \homebuttons
\end{frame}
\begin{frame}
    \frametitle{Usage}
    \hypertarget<1->{Usage}{}
    \begin{itemize}
        \item Use a suitable PDF viewer in presentation mode
            \begin{itemize}
                \item Buttons may not work in all PDF viewers
                \item Tested works with: Evince
                \item Tested doesn't work with: Chromium
                \item (buttons can be quite slow-- try not to spam them)
            \end{itemize}
        \item There is an A side and B side (like a vinyl)
        \item Navigate by clicking the buttons
    \end{itemize}
    \homebuttons
\end{frame}
\section{A}
{
    \usebackground{horimiya-masks.jpg}
    % AHome
    \begin{frame}
        \hypertarget<1->{AHome}{}
        \begin{tikzpicture}[remember picture, overlay]
            \node[shift={(5.7cm, 6.5cm)}]() at (current page.south west){
                \hyperlink{masks}{\beamerbutton{masks}}
            };
            \node[shift={(9.8cm, 3.8cm)}]() at (current page.south west){
                \hyperlink{traditionalClothing}{\beamerbutton{traditional clothing}}
            };
            \node[shift={(6.9cm, 8.8cm)}]() at (current page.south west){
                \hyperlink{su}{\beamerbutton{sailor uniform}}
            };
            \node[shift={(6.5cm, 1.4cm)}]() at (current page.south west){
                \hyperlink{sw}{\beamerbutton{street wear}}
            };
            \node[shift={(7.5cm, 1.8cm)}]() at (current page.south west){
                \hyperlink{horimiya}{\beamerbutton{horimiya}}
            };
        \end{tikzpicture}
        \homebuttons
    \end{frame}
}
{
% masks
    \usebackground[0.3]{horimiya-masks.jpg}
    \begin{frame}
        \hypertarget<1->{masks}{}
        \begin{tikzpicture}[remember picture, overlay]
            \node[draw, align=center, fill=blue, text=white, shift={(5.7cm, 6.5cm)}]() at (current page.south west){
                    On this volume cover, the eponymous\\ characters wear festival masks
            };
            \node[shift={(9.5cm, 4.5cm)}]() at (current page.south west){
                \hyperlink{masksWhat}{\beamerbutton{what are they?}}
            };
            \node[shift={(9.5cm, 2.5cm)}]() at (current page.south west){
                \hyperlink{masksCover}{\beamerbutton{why on a volume cover?}}
            };
        \end{tikzpicture}
        \homebuttons
    \end{frame}
}
{
% masksWhat
    \usebackground[0.3]{horimiya-masks.jpg}
    \begin{frame}
        \hypertarget<1->{masksWhat}{}
        \begin{tikzpicture}[remember picture, overlay]
            \node[draw, align=center, fill=blue, text=white, shift={(9.5cm, 4.5cm)}]() at (current page.south west){
                    These masks are often worn during\\ festivals and shrine events (\cite{ct}).\\
                    They come from folklore or Noh theater.
            };
            \node[shift={(-2.0cm, 0.5cm)}]() at (current page.south east){
                \hyperlink{masks}{\beamerbutton{Back}}
            };
        \end{tikzpicture}
        \homebuttons
    \end{frame}
}
{
% masksCover
    \usebackground[0.3]{horimiya-masks.jpg}
    \begin{frame}
        \hypertarget<1->{masksCover}{}
        \begin{tikzpicture}[remember picture, overlay]
            \node[draw, align=center, fill=blue, text=white, shift={(9.5cm, 2.5cm)}]() at (current page.south west){
                    Paired with the kimonos,\\
                    the festival masks elicit a sense of excitedness and enjoyment.\\
                    It's an escape from their usual school uniform-clad days.
            };
            \node[shift={(-2.0cm, 0.5cm)}]() at (current page.south east){
                \hyperlink{masks}{\beamerbutton{Back}}
            };
        \end{tikzpicture}
        \homebuttons
    \end{frame}
}
{
% traditional clothing
    \usebackground[0.3]{horimiya-masks.jpg}
    \begin{frame}
        \hypertarget<1->{traditionalClothing}{}
        \begin{tikzpicture}[remember picture, overlay]
            \node[draw, align=center, fill=blue, text=white, shift={(9.8cm, 3.8cm)}]() at (current page.south west){
                On this volume cover, the eponymous\\ characters wear traditional kimonos
            };
            \node[shift={(6.0cm, 7.5cm)}]() at (current page.south west){
                \hyperlink{kimonoRole}{\beamerbutton{role of kimonos today}}
            };
            \node[shift={(6.0cm, 5.5cm)}]() at (current page.south west){
                \hyperlink{kimonoStyle}{\beamerbutton{kimono styles}}
            };
            \node[shift={(6.0cm, 2.0cm)}]() at (current page.south west){
                \hyperlink{kimonoCover}{\beamerbutton{why on a volume cover?}}
            };
        \end{tikzpicture}
        \homebuttons
    \end{frame}
}
{
% kimonoRole
    \usebackground[0.3]{horimiya-masks.jpg}
    \begin{frame}
        \hypertarget<1->{kimonoRole}{}
        \begin{tikzpicture}[remember picture, overlay]
            \node[draw, align=center, fill=blue, text=white, shift={(6.0cm, 5.5cm)}]() at (current page.south west){
                Traditionally, \textit{kimono} encompassed both\\ everyday wear and more special attires,\\
                but the modern kimono solely comes from\\ the decorative festive wear of the\\ premodern elite
                (\cite{gg} p.353).\\
                Nowadays, kimono are worn primarily for formal special occasions.\\
                Modern kimono have come not just to be a\\ symbol of Japaneseness,
                but also of feminine etiquette,\\ socialization, and conformance for women (\cite{gg} p.362).\\
                For the annual coming-of-age ceremony,\\
                twenty year old women are expected to wear elaborate kimonos,\\
                while men are expected to wear a Western suit (\cite{gg} p.352).
            };
            \node[shift={(-2.0cm, 0.5cm)}]() at (current page.south east){
                \hyperlink{traditionalClothing}{\beamerbutton{Back}}
            };
        \end{tikzpicture}
        \homebuttons
    \end{frame}
}
{
% kimonoStyle
    \usebackground[0.3]{horimiya-masks.jpg}
    \begin{frame}
        \hypertarget<1->{kimonoStyle}{}
        \begin{tikzpicture}[remember picture, overlay]
            \node[draw, align=center, fill=blue, text=white, shift={(6.0cm, 5.5cm)}]() at (current page.south west){
                Kimono vary depending on their purpose.\\
                They can be made of silk, cotton, or linen.\\
                Patterns invoke seasons or certain artistic tastes (\cite{patternBook}).\\
                Notably, floral patterns are reserved for the refined woman
                (\cite{jDaily}).\\
                Interestingly, Kyoko (the girl on the cover) wears a plain kimono\\
                with tasuki, cloth tied around the sleeves to improve mobility.\\
                Compared to having her wear a floral restricting kimono,\\
                this simpler more active style better represents her character.
            };
            \node[shift={(-2.0cm, 0.5cm)}]() at (current page.south east){
                \hyperlink{traditionalClothing}{\beamerbutton{Back}}
            };
        \end{tikzpicture}
        \homebuttons
    \end{frame}
}
{
% kimonoCover
    \usebackground[0.3]{horimiya-masks.jpg}
    \begin{frame}
        \hypertarget<1->{kimonoCover}{}
        \begin{tikzpicture}[remember picture, overlay]
            \node[draw, align=center, fill=blue, text=white, shift={(6.0cm, 2.0cm)}]() at (current page.south west){
                Manga artists like to draw their characters in different outfits and\\
                especially ones the characters don't typically wear.\\
                But more than that, as kimonos today are\\ associated with festivals,
                this cover radiates high energy--\\
                an attractive quality for this genre of light fiction.
            };
            \node[shift={(-2.0cm, 0.5cm)}]() at (current page.south east){
                \hyperlink{traditionalClothing}{\beamerbutton{Back}}
            };
        \end{tikzpicture}
        \homebuttons
    \end{frame}
}
{
% su
    \usebackground[0.3]{horimiya-masks.jpg}
    \begin{frame}
        \hypertarget<1->{su}{}
        \begin{tikzpicture}[remember picture, overlay]
            \node[draw, align=center, fill=blue, text=white, shift={(6.9cm, 6.8cm)}]() at (current page.south west){
                Uniforms are pervasive in Japan.\\
                Employee uniforms and even\\ job-hunting uniforms are common.\\
                School uniforms are no exception:\\
                Elementary, middle, and high schools\\ each have their own typical uniforms.\\
                The sailor uniform is common for middle schools
                \cite{nippon}.
            };
        \end{tikzpicture}
        \homebuttons
    \end{frame}
}
{
% sw
    \usebackground[0.3]{horimiya-masks.jpg}
    \begin{frame}
        \hypertarget<1->{sw}{}
        \begin{tikzpicture}[remember picture, overlay]
            \node[draw, align=center, fill=blue, text=white, shift={(6.5cm, 1.4cm)}]() at (current page.south west){
                The main pair donning street wear.\\
                In contrast with the main cover,\\
                the trench coat and pants,\\ and jacket and skirt are Western.
            };
        \end{tikzpicture}
        \homebuttons
    \end{frame}
}
{
% horimiya
    \usebackground[0.3]{horimiya-masks.jpg}
    \begin{frame}
        \hypertarget<1->{horimiya}{}
        \begin{tikzpicture}[remember picture, overlay]
            \node[draw, align=center, fill=blue, text=white, shift={(7.5cm, 1.8cm)}]() at (current page.south west){
                    Horimiya is a romantic comedy\\ manga written/set in late 2000's Tokyo.\\
                    The main pair are in their last year of high school.\\
                    The title is a concatonation of\\ the names of the main characters:\\
                    Hori (left) and Miyamura (right).
            };
        \end{tikzpicture}
        \homebuttons
    \end{frame}
}
\section{B}
{
    \usebackground{horimiya-school.jpg}
    % BHome
    \begin{frame}
        \hypertarget<1->{BHome}{}
        \begin{tikzpicture}[remember picture, overlay]
            \node[shift={(10.5cm, 3.4cm)}]() at (current page.south west){
                \hyperlink{usage}{\beamerbutton{overwork}}
            };
            \node[shift={(6.4cm, 2.8cm)}]() at (current page.south west){
                \hyperlink{usage}{\beamerbutton{school uniforms}}
            };
        \end{tikzpicture}
        \homebuttons
    \end{frame}
}
{
% overwork
    \usebackground[0.2]{horimiya-school.jpg}
    \begin{frame}
        \homebuttons
    \end{frame}
}
\section{Sources}
\begin{frame}
    \frametitle{Sources}
    \begin{thebibliography}{00}
        \setbeamertemplate{bibliography item}{\(\blacktriangleright\)}
        \bibitem{horimiyaManga}
        Horimiya manga
        \begin{itemize}
            \item A: Volume 16, Chapter 120, Page 1 (Volume Cover)
            \item B: Volume 13, Chapter 100, Page 2
        \end{itemize}
        \bibitem{gg}
        Ofra Goldstein-Gidoni.
        \textit{Kimono and the Construction of Gendered and Cultural Identities}.
        Ethnology, Autumn, 1999, Vol. 38, No. 4 (Autumn, 1999), pp. 351-370
        \bibitem{jDaily}
        Cynthia Green,
        \href{https://daily.jstor.org/the-surprising-history-of-the-kimono/}{https://daily.jstor.org/the-surprising-history-of-the-kimono}
        \bibitem{patternBook}
        Betty Y. Siffert.
        \textit{Hinagata Bon: The Art Institute of Chicago Collection of Kimono Pattern Books}.
        Art Institute of Chicago Museum Studies , 1992, Vol. 18, No. 1, Five Centuries of Japanese Kimono: On This Sleeve of Fondest Dreams (1992), pp. 86-94+103
        \bibitem{ct}
        Alicia Joy,
        \href{https://theculturetrip.com/asia/japan/articles/traditional-japanese-masks-and-what-theyre-used-for/}{https://theculturetrip.com/asia/japan/articles/traditional-japanese-masks-and-what-theyre-used-for}
        \bibitem{nippon}
        Mitamura Fukiko,
        \href{https://www.nippon.com/en/currents/d00216/}{https://www.nippon.com/en/currents/d00216}
    \end{thebibliography}
\end{frame}
\end{document}
